\section{Motivation}
\label{sec:intro-motivation}

Consider a change that narrows what a widely used helper function will accept:
where it previously tolerated a missing value, it now requires one. Inside the
diff the change is sound. The function's own tests are updated, and a reviewer
reading those few lines has every reason to approve. The defect, if there is
one, is somewhere else: in a caller three directories away that has always
passed nothing and relied on the old tolerance. Nothing in the diff mentions
that caller.

This scenario illustrates a recognised challenge in code review; its frequency
is not estimated in this thesis. When Bacchelli and Bird asked Microsoft
developers what they hoped to get from review and what they actually got, the
gap was defect detection, and the
reason they gave was not inattention but
\emph{understanding}~\cite{bacchelli2013}: a reviewer who does not know what
depends on a piece of code cannot judge whether changing it is safe.

A language model asked to review a pull request is usually given the diff and
nothing else. It cannot open the neighbouring file, search for usages, or run
the tests. The natural response is to put more of the repository into the
prompt, and two mechanisms are available. Retrieval embeds the repository and
returns the code most similar to the change. Structural traversal parses the
repository and returns the code \emph{connected} to the change. In the
vignette above, similarity would rank the helper's sibling utilities highly,
while the caller that is about to break may share no vocabulary with the
function it calls. Which mechanism to use, and what each is actually good for,
is the question this thesis takes up.

\section{Problem Statement}
\label{sec:intro-problem}

The problem is not that repository context is unavailable. Both retrieval and
graph construction are mature enough to deploy, and systems built on them have
demonstrated gains on completion, fault localisation and
repair~\cite{zhang2023repocoder,liu2024graphcoder,yang2025kgcompass}. The
problem is that for review-comment generation it is not known which kind of
context helps, by how much, or whether the instruments commonly used to answer
such questions can detect the answer.

Three difficulties compound. First, the tasks on which structural context has
been shown to help all have checkable outcomes (completion is right or
wrong, a test passes or fails), while a review comment is prose with no
automatic success signal, making evaluation part of the problem. Second, the
evaluation methods now used for this task are weaker than the claims made with
them: single-model rubric scoring with unreported reliability cannot
distinguish a real improvement from a preference for longer
text~\cite{zheng2023judging}. Third, the structural context must come from
a described static-analysis artefact, here a code property graph, rather
than an ad hoc text heuristic, so that the result can be attributed and
replicated.

\section{Research Questions}
\label{sec:intro-rqs}

Three research questions form a chain, each closing a gap the previous one
leaves open:

\begin{description}
    \item[RQ1: Comparative coverage.] \textit{How do \kg{}, \rag{}, and
        \hybrid{} context affect the rubric coverage of LLM-generated review
        comments compared with a diff-only \baseline{}?}
    \item[RQ2: Cross-file reasoning.] \textit{To what extent does
        structural context enable an LLM reviewer to identify cross-file
        consequences that are absent from the diff?}
    \item[RQ3: Human preference.] \textit{Do human raters prefer
        graph-augmented reviews over diff-only reviews for identifying affected
        code when the review conditions are hidden?}
\end{description}

RQ1 is evaluated with a 25-criterion rubric, including nine criteria
designated in advance as structurally relevant (Section~\ref{sec:rubric}).
The rubric measures what a review covers, not whether it is correct, so
RQ2 uses controlled defect injection, where the ground truth is known.
Both are scored by language models, so RQ3 asks people. RQ1 and RQ2 build the
knowledge graph with Joern; RQ3 uses the same graph block with edges
resolved from import statements. The hybrid arm tests whether structural and similarity context
combine their benefits or dilute attention when concatenated.

\section{Contributions}
\label{sec:intro-contributions}

\paragraph{Mechanistic evidence for the structural blind spot.}
A controlled injection experiment shows that a diff-only reviewer detects none
of 28 cross-file defects, while the graph arm detects 18. Near-ceiling
performance on local controls supports a structural interpretation of the gain.

\paragraph{A localised benefit on pull requests.}
On 35~pull requests, the graph arm concentrates its gain in
the nine criteria its mechanism predicts, while retrieval produces the largest
full-rubric gain. The hybrid mode leads on neither scale after correction.

\paragraph{Human corroboration of targeted grounding.}
In the human study, raters prefer the graph-augmented review for naming
affected code, but show no preference on overall usefulness; code clarity
favours the baseline.

\paragraph{An analysis of evaluation sensitivity.}
Judge-level and regeneration analyses show that the targeted direction is
stable, but its magnitude and significance depend on the evaluation
configuration.

\medskip\noindent
The remainder of this thesis is organised as follows.
Chapter~\ref{chap:background} situates the work in the literature on code
review, language models for code, retrieval, knowledge graphs, and evaluation.
Chapter~\ref{chap:approach} describes the pipeline and the graph
construction. Chapter~\ref{chap:edas} presents the experimental design, the
rubric, and the human study protocol. Chapter~\ref{chap:results} reports and
discusses results from all three studies together with threats to validity.
Chapter~\ref{chap:conclusion} summarises the findings and identifies future
work.
